\chapter{线宽、相干性、光束质量与亮度}
\label{ch:metrics}

\section*{学习目标}
\begin{itemize}
  \item 从严格定义出发理解线宽、相干性、$M^2$、亮度、偏振纯度与工作点之间的关系。
  \item 明确哪些性能量可以直接由确定性电磁/器件求解器给出，哪些量只能借助降阶噪声模型间接估计。
  \item 学会把性能指标写成“条件化、可复核、可对比”的结果，而不是脱离工作点的孤立数字。
\end{itemize}

\section*{先修知识}
建议已掌握 \cref{ch:threshold,ch:single-mode-design,ch:robustness,ch:dynamics,ch:eto-coupling} 中关于阈值、单模工作窗口、鲁棒性窗口、动态稳定性和热-电-光反馈的内容。

\section*{本章路线图}
前面两章已经分别回答了“名义设计对工艺扰动是否稳健”以及“起振以上的动态工作窗口如何理解”，而本章要回答的是“这些内部动态最终怎样投影成外界真正测到的指标”。这一步很容易再次发生层级混乱，因为实验测到的是线宽、远场、偏振、功率和亮度，而模型里最先算出的往往是复频率、$Q$、模场、温度场以及上一章中的动态中间量。我们将先建立性能指标的严格定义，再说明这些定义如何与 PCSEL 的工作点连接，随后专门讨论仿真和实验之间的接口，最后回答一个实践中最常被问到的问题：哪些量可以在 Lumerical 和 COMSOL 中直接算，哪些量只能通过它们提供的中间量来间接估计。

\section{为什么性能指标必须和工作点一起报告}
对于普通课堂中的理想谐振腔，给出一个本征频率和一个 $Q$ 往往已经能说明很多问题。对 PCSEL 来说，这远远不够。因为真实器件工作在外延层、电流路径和热反馈共同参与的环境里，同一个结构在不同电流、不同散热条件、不同封装反射和不同观测时间窗下，线宽、偏振纯度和远场都可能明显改变。

这意味着性能指标不是脱离条件存在的“器件属性标签”，而是工作状态的函数。若不先说明工作点，就谈“该器件的线宽是多少”“该器件的光束质量是多少”，这种说法通常不严谨。教材式写法应当始终采用条件化表述，例如：“在 \SI{25}{\celsius}、连续波、某电流区间、某观测窗口下，线宽与偏振纯度落在什么范围内。”

这个原则之所以重要，不只是为了报告规范，也为了防止错误建模。因为只要读者接受了“性能是工作点的函数”这一点，后面就不会轻易把冷腔结果直接冒充工作状态下的器件结论。

\begin{IntuitionBox}
对 PCSEL 来说，性能指标不是一张独立成绩单，而是对“模式竞争、注入分布和热反馈当前处于什么状态”的可观测投影。看指标时必须同时问：它是在什么条件下测到或算到的。
\end{IntuitionBox}

\section{从一阶相干函数到线宽：定义先于经验公式}
设输出复电场为 $E(t)$，一阶相干函数定义为
\begin{equation}
 g^{(1)}(\tau)=\frac{\langle E^*(t)E(t+\tau)\rangle}{\langle |E(t)|^2\rangle}.
\end{equation}
若在某个近似区间内相位扩散可看作平稳随机过程，那么频谱与 $g^{(1)}(\tau)$ 互为傅里叶变换。在线形接近 Lorentz 形时，常把谱半高全宽记为线宽 $\Delta\nu$。这个定义看似标准，真正困难的是：PCSEL 常常同时受到快噪声和慢漂移影响，因此“线宽”可能依赖观测时间窗和拟合方法。

若进一步采用最经典的相位扩散模型，
\begin{equation}
g^{(1)}(\tau)=\exp(-\pi\Delta\nu_{\mathrm{L}}|\tau|),
\label{eq:g1_lorentz}
\end{equation}
其傅里叶变换就是 Lorentz 谱线，而相干时间满足
\begin{equation}
\tau_c \approx \frac{1}{\pi\Delta\nu_{\mathrm{L}}}.
\label{eq:coherence_time_lorentz}
\end{equation}
\cref{eq:coherence_time_lorentz} 只在 Lorentz 线形、单模、近稳态相位扩散主导的条件下成立；一旦慢漂移主导，$\tau_c$ 与“某个单独线宽数字”之间就不再能用这一条关系简单换算。

为了避免概念跳跃，最好把观测线宽写成分层诊断量
\begin{equation}
\Delta\nu_{\mathrm{obs}}
\approx
\Delta\nu_{\mathrm{int}}
+
\Delta\nu_{\mathrm{thermal}}
+
\Delta\nu_{\mathrm{env}},
\label{eq:linewidth_layers}
\end{equation}
其中 $\Delta\nu_{\mathrm{int}}$ 表示内禀相位噪声主导的部分，$\Delta\nu_{\mathrm{thermal}}$ 表示温度起伏和慢热漂移引起的展宽，$\Delta\nu_{\mathrm{env}}$ 则是外腔反射、机械振动和电源噪声等外禀项。这种写法不是严格可加定理，而是\textbf{实验与建模对接时的诊断分解}。它的意义在于提醒读者：只用一个瞬时测得的线宽数字，往往无法回答性能差异究竟来自腔本身、热反馈还是测试环境。

在半导体激光器近似理论中，人们常写出与输出功率、光子寿命和线宽增强因子有关的估计式。这里最常用的线宽增强因子就是 Henry 因子
\begin{equation}
\alpha_H = -2k_0\frac{\partial n/\partial N}{\partial g/\partial N},
\label{eq:henry_factor_metrics}
\end{equation}
其中 $k_0=\omega/c$。对 GaAs 量子阱，$\alpha_H$ 常落在 $2$--$5$ 的量级，因此实际线宽往往会比理想 Schawlow--Townes 极限再乘上一个 $(1+\alpha_H^2)$ 级别的放大因子。这个估计有助于建立量纲感和趋势判断，但在 PCSEL 中不能被机械套用，原因有两点。第一，$\alpha_H$ 的标准推导基于单模行波激光器中振幅扰动通过载流子变化回写到相位这一图像；PCSEL 工作在驻波/区中心模式，相位-幅度耦合的空间结构与行波激光器不同，标准单模 $\alpha_H$ 不能直接描述多波耦合模的相位响应。第二，PCSEL 多模竞争、偏振切换和热漂移会引入额外的低频相位调制，这些贡献不包含在 $\alpha_H$ 的定义中。因此，一旦工作点接近偏振切换、模式跳变或强热漂移，以单模平稳近似为基础的 $(1+\alpha_H^2)$ 线宽放大公式就会失效。

\section{频率噪声谱、短时/长时线宽与测量协议}
若只报告一个“线宽”数字，而不说明它来自哪一种测量时间窗和拟合约定，结果往往不可比。更稳妥的现代写法，是先看频率噪声功率谱密度 $S_\nu(f)$，再讨论在给定观测协议下它如何投影成谱线宽。定性上说，高频白频率噪声更容易对应 Lorentz 型展宽，低频漂移和慢热起伏则更容易把谱线推向 Gaussian 或 Voigt 型轮廓。于是，“短时线宽”和“长时线宽”在实验上往往并不是同一个数字。

把这条链写成最小公式，就是先引入相位与瞬时频率偏移的关系
\begin{equation}
\delta \nu(t)=\frac{1}{2\pi}\frac{\dd \phi}{\dd t},
\label{eq:freq_phase_relation}
\end{equation}
再用频率噪声谱 $S_\nu(f)$ 区分快慢贡献。若高频端可近似为白频率噪声，且采用常用单边谱约定，则对应的 Lorentz 线宽可写成
\begin{equation}
\Delta\nu_{\mathrm{L}} \approx \pi S_\nu^{(\mathrm{white})}.
\label{eq:lorentz_from_freqnoise}
\end{equation}
若低频端主要由慢漂移决定，频率分布更接近高斯，则常把 Gaussian 半高全宽写成
\begin{equation}
\Delta\nu_{\mathrm{G}} = 2\sqrt{2\ln 2}\,\sigma_\nu,
\label{eq:gaussian_linewidth}
\end{equation}
其中 $\sigma_\nu$ 是观测窗口内的频率标准差。于是，当白噪声与慢漂移同时存在时，实验上常把谱线视为 Voigt 轮廓，并用近似公式
\begin{equation}
\Delta\nu_{\mathrm{V}}
\approx
0.5346\,\Delta\nu_{\mathrm{L}}
+
\sqrt{0.2166\,\Delta\nu_{\mathrm{L}}^2+\Delta\nu_{\mathrm{G}}^2}
\label{eq:voigt_width}
\end{equation}
估计总半高全宽。\cref{eq:lorentz_from_freqnoise,eq:gaussian_linewidth,eq:voigt_width} 的真正价值，不是要求读者把数字系数背下来，而是帮助读者建立一条清楚的诊断链：\textbf{先看频率噪声来自快白噪声还是慢漂移，再决定线宽应按 Lorentz、Gaussian 还是 Voigt 去解释。}

对教材读者而言，最小可执行协议至少应报告四件事：观测时间窗、频谱分辨率带宽、拟合线形（Lorentz / Gaussian / Voigt），以及是否做了偏振选择和模式选择。只有在这四个条件明确后，线宽比较才有物理意义。否则，同一器件在不同仪器设置下得到两个明显不同的“线宽”，并不一定意味着器件状态真的变了。

\begin{RigorousBox}
频率噪声谱与线宽之间不是“一条公式永远通吃”的关系。更合理的理解是：线宽是测量协议下对频率噪声过程的压缩表征。因此，只给线宽而不给时间窗、分辨率和拟合约定，是不完整的实验报告。
\end{RigorousBox}

\section{Petermann 因子与开放腔非正交性：什么时候必须提醒读者}
对开放腔激光器，尤其是具有明显非厄米性质的模式，线宽讨论还会遇到一个额外修正语言：Petermann 因子。它最小地回答一个问题：\textbf{当激光模并非理想正交模时，噪声怎样被额外放大。} 对 PCSEL 而言，这个语言不是每次都要展开成主公式，但至少要在教材里说明它的适用语境。

更稳妥的表述是：若目标模表现出明显开放性、非正交性增强，或器件工作点接近强非厄米模混合区域，那么把线宽只写成“Schawlow--Townes 乘一个 Henry 修正”可能不够。此时 Petermann 因子提供了一种提醒：开放腔中的模非正交会让相位噪声放大。它并不是大面积 PCSEL 的默认主控量，但在讨论高 $Q$、强开放、近简并或强模式混合区域时，不应完全缺席。

\section{外腔反馈与环境耦合：为什么性能章也必须给出一句硬说明}
\cref{ch:dynamics} 已经把外腔反馈放进动态扰动入口，但本章还应从性能报告角度再补一句更硬的话：\textbf{任何线宽、相干性或稳定性指标，若未说明外腔反射与封装环境，就不应被自动读成“器件内禀极限”。}

这是因为 PCSEL 的面发射几何天然更容易和自由空间光路、窗口片和封装界面发生回馈耦合。实验上经常出现的现象是：冷腔和工作点都没有明显变化，但一旦更换外部光路或测试台，低频频率噪声、长时线宽和偏振稳定性就一起改变。对教材读者来说，最重要的不是在这里学会完整外腔理论，而是建立报告纪律：\textbf{内禀量、外禀量、短时量和长时量必须分开。}

\begin{RigorousBox}
线宽是\textbf{谱形定义}，不是“任何频谱图上随手量一个宽度”。若没有说明观测窗口、采样时间、频谱分辨率和拟合线形，就不能把不同文献或不同实验条件下的线宽直接拿来比较。
\end{RigorousBox}

\section{相干性与线宽相关，但并不完全等价}
教材中常把“线宽越窄，相干性越好”写成一句简单结论。这个说法在理想单模、稳定噪声主导的场景里有一定道理，但在 PCSEL 里需要补足两个限定。

第一，相干性有时间相干和空间相干两个侧面。线宽主要对应时间相干，而大面积 PCSEL 的空间相干则与横向模分布、相位一致性和模式竞争更直接相关。一个器件即使在短时间窗内线宽不宽，若横向上已经出现多模分区或局部失相干，其远场与亮度仍然会明显变差。

第二，线宽本身也可能同时混入慢漂移和快噪声。若观测窗口从微秒扩展到秒级，谱宽可能显著增加，但这不一定意味着腔本身的内禀噪声变大，而可能是温控和封装反射在更长时间尺度上显现出来。

因此，在 PCSEL 的研究报告中，最合理的做法不是孤立汇报“线宽”两个字，而是同时给出观测时间窗、偏振条件、工作电流和温度。只有这样，线宽才真正能与相干性和模式稳定性联系起来。

\section{\texorpdfstring{$M^2$}{M-squared}、亮度和功率为什么必须一起看}
高功率并不自动等于高亮度，这是大面积面发射器件最容易出现的误判。亮度常写为
\begin{equation}
B = \frac{P}{A\Omega},
\label{eq:brightness_def}
\end{equation}
其中 $P$ 是输出功率，$A$ 是有效出光面积，$\Omega$ 是发散对应的立体角。若器件面积增大导致输出功率上升，但横向多模使远场发散明显变宽，则亮度可能不升反降。于是，对 PCSEL 来说，功率、$M^2$ 和亮度必须一起看，而不能只报其中一个最好看的数字。

$M^2$ 则提供了把真实光束与理想高斯基模比较的桥梁。对于大面积 PCSEL，它是空间相干和模式纯度的一个非常直观的外部指标。若同一器件在电流升高后 $M^2$ 明显变坏，同时偏振纯度下降、远场主瓣漂移增强，那么最合理的解释通常不是“输出只是更强了”，而是模式竞争和热反馈已经开始破坏目标模的主导地位。

若采用二阶矩定义，并在某一横向方向上记近场标准差为 $\sigma_x$、远场角谱标准差为 $\sigma_{\theta_x}$，则
\begin{equation}
M_x^2=\frac{4\pi}{\lambda}\sigma_x\sigma_{\theta_x},
\qquad
M_y^2=\frac{4\pi}{\lambda}\sigma_y\sigma_{\theta_y}.
\label{eq:m2_second_moment}
\end{equation}
对理想高斯基模有 $M_x^2=M_y^2=1$。若光束近似圆对称、且希望做快速数量级判断，也常写成
\begin{equation}
M^2\approx \frac{\pi w_0\theta}{\lambda},
\label{eq:m2_quick_est}
\end{equation}
其中 $w_0$ 是近场束腰半径，$\theta$ 是远场半角。\cref{eq:m2_second_moment,eq:m2_quick_est} 的区别应在教材里说清：前者更严格，适合实验报告与复杂光束；后者更直观，适合教学估算。

\begin{ExampleBox}
若某器件从 \SI{1}{A} 升到 \SI{2}{A} 时总功率上升 50\%，但远场半角扩大一倍、$M^2$ 明显变差，那么亮度通常不会同步上升。此时最重要的不是继续追求更高电流，而是检查目标模是否已经失去对竞争模的净优势。
\end{ExampleBox}

\begin{ExampleBox}
再给一个最小数量级例子。若把大面积 PCSEL 的主出光近场近似为直径 $D$ 的均匀圆孔，则衍射极限半角可粗略写成
\[
\theta_{1/e^2}\sim \frac{2\lambda}{\pi D}.
\]
取 $\lambda=\SI{980}{nm}$、$D=\SI{200}{\micro\meter}$，得到
\[
\theta_{1/e^2}\sim 3\times10^{-3}\ \mathrm{rad}\approx 0.18^\circ.
\]
这个数字的作用不是给所有器件套统一标准，而是提供一个基线：若实测发散角远大于该量级，就不应先把原因归咎于“大面积”，而应优先检查近场幅度不均、相位起伏、边缘泄漏、热透镜或多模共存。
\end{ExampleBox}

\section{偏振纯度、远场稳定性和模式稳定性本质上连在一起}
PCSEL 的偏振与远场并不是两个孤立测量量。它们都来自近场对称性和辐射通道的组织方式，而这又会被注入和温升动态改写。于是，在工作点扫描中常见到这样的联动：偏振纯度下降的同时，远场主瓣变宽或偏移；又或者远场出现副瓣增强的同时，线宽也开始展宽。这些现象不应被当作互不相干的“多个问题”，它们通常是同一个模式竞争过程在不同观测端口上的投影。

因此，本章讨论性能时必须和前一章连起来看：当目标模相对于竞争模的裕量缩小时，最先暴露异常的，不一定是阈值或总功率，而往往是偏振纯度、远场或短时间稳定性。性能指标之所以是高价值指标，恰恰因为它们能比“器件还能不能亮”更早地暴露设计边界。

\section[直接与间接可计算量]{哪些量可以在 Lumerical 和 COMSOL 中直接算，哪些只能间接估计}
这一节直接回答建模实践中的常见问题。Lumerical 和 COMSOL 本质上都以确定性电磁场、输运方程和热传导方程为核心。它们最擅长的，是给出\textbf{模式、损耗、场分布、载流子分布、温度分布以及这些量对参数的灵敏度}。而线宽和相干性则牵涉随机噪声与相位扩散，通常不能直接通过一次确定性全波仿真“求出来”。

\cref{tab:metric_toolmap} 给出一张实用区分表。

\begin{table}[htbp]
  \centering
  \footnotesize
  \caption{性能指标与 Lumerical / COMSOL 的典型建模接口。重点在于区分“可直接计算量”和“需降阶估计量”。}
  \label{tab:metric_toolmap}
  \setlength{\tabcolsep}{4pt}
  \begin{tabularx}{\textwidth}{@{}p{1.8cm}p{2.1cm}Y Y@{}}
    \toprule
    指标 & 物理来源 & Lumerical 中常见做法 & COMSOL 中常见做法 \\
    \midrule
    共振频率、模场、$Q$ & 冷腔或含损耗光学 & FDTD / \software{MODE} / RCWA 提取复频率、衰减率、场分布 & \software{Wave Optics} 本征频率或频域求解 + PML \\
    远场、偏振、上下出光比例 & 开放辐射通道 & 周期 RCWA 或有限阵列 FDTD + near-to-far & \software{Wave Optics} 频域 + 端口 / 远场后处理 \\
    约束因子、损耗重叠 & 模场与层栈重叠 & \software{MODE} / FDTD 后处理重叠积分 & \software{Wave Optics} 场积分后处理 \\
    电流扩展、J-V、载流子分布 & 漂移-扩散电学 & \software{CHARGE} & \software{Semiconductor} \\
    温度场、热漂移灵敏度 & 热传导 + 热源 & \software{HEAT} & \software{Heat Transfer} \\
    线宽、相干时间 & 噪声与相位扩散 & 由 $Q$、功率、热灵敏度等中间量进入降阶噪声模型 & 同左，通过中间量进入降阶模型 \\
    亮度、$M^2$ & 远场与功率联合定义 & 由远场与功率后处理得到 & 由远场与功率后处理得到 \\
    \bottomrule
  \end{tabularx}
\end{table}

如果要说得再直接一点：
\begin{itemize}
  \item 想算候选模、远场、偏振、被动损耗和有限阵列出射，优先使用光学求解器；
  \item 想算电流扩展、阈值电流的器件级来源、串联电阻和非均匀增益，必须进入 \software{CHARGE} 或 \software{Semiconductor} 一类漂移-扩散求解器；
  \item 想算温升、热源和热漂移，必须进入 \software{HEAT} 或 \software{Heat Transfer} 一类热求解器；
  \item 想直接“从一次仿真得到线宽”，通常做不到，只能用上述求解器提供的中间量建立降阶噪声模型。
\end{itemize}

这一点对初学者尤其关键，因为它能立刻避免一个常见误区：把确定性全波仿真结果误读成完整的相干性结论。严格地说，全波求解器给出的是线宽模型所需的参数，而不是线宽本身的全部统计信息。

\section{性能报告的正确写法：条件、区间和可复核性}
到了收束阶段，最值得坚持的写法是把结果写成“条件 + 区间 + 提取方法”的结构。与其写“器件线宽为多少、亮度为多少”，不如写成：在何种电流、温度、观测窗口和偏振选择下，线宽、$M^2$、偏振纯度和远场主瓣角分别落在哪些范围，提取算法是什么，仿真或实验是否重复验证过。

这种写法看起来更繁琐，但它有两个长远好处。第一，它天然阻止跨层级偷换概念，因为每个指标都带着自己的求解或测量上下文。第二，它使模拟与实验更容易对接，因为双方都知道哪些中间条件必须一致才能比较。

\begin{PitfallBox}
把某一电流点、某一时间窗、某一偏振状态下得到的“最优线宽”当成器件代表性能，是典型的过度报告。对于 PCSEL，这种写法几乎必然掩盖模式竞争和热漂移带来的窗口效应。
\end{PitfallBox}

\section*{本章小结}
PCSEL 的性能指标必须在工作点上下文中理解。线宽来自谱形定义并依赖观测窗口，相干性既有时间侧也有空间侧，$M^2$ 和亮度要求把功率与发散一起看，偏振纯度和远场稳定性则是模式稳定性的外部投影。频率噪声谱、短时/长时线宽、Petermann 因子和外腔反馈进一步提醒我们：开放腔激光器的“性能数字”总是带着测量协议和环境边界条件。在建模层面，Lumerical 与 COMSOL 都擅长直接计算模式、辐射、载流子和温度等确定性中间量，但线宽和相干性通常需要借助这些中间量进入降阶噪声模型间接估计。真正可靠的性能报告，不是“报一个最好看的数字”，而是报告一个可复核的工作窗口。

\section*{练习题}
\begin{enumerate}
  \item \basicq 解释为什么在 PCSEL 中，线宽不能脱离观测时间窗单独比较。

  \item \basicq 结合 \cref{eq:brightness_def} 说明为什么高功率并不必然意味着高亮度。

  \item \midq 试着把“偏振纯度下降、远场变宽、线宽增加”三件事串成一条可能的模式竞争链。

  \item \hardq 选择一个你关心的性能指标，说明它在 Lumerical 或 COMSOL 中是直接计算量，还是需要通过中间量间接估计。
\end{enumerate}

\section*{建议继续阅读}
\begin{itemize}
  \item 下一章将把前面所有理论主线压缩成一套真正可执行的 PCSEL 仿真工作流，并明确不同软件链如何协同。 这条阅读的重点是把本章的输出顺着主线接到后续模块，而不是停成孤立知识点。

  \item 若想更深入理解半导体激光器中的线宽与相干，可参考 \cite{siegman1986,coldren2012,henry1982} 中关于相干函数、幅相耦合、线宽和光束质量的内容。 这条阅读的重点是补足本章关于线宽、相干性、亮度与报告规范 的标准背景、经典语言和代表性文献接口。
\end{itemize}
